\begin{table}[htbp]
\centering
\small
\caption{Mean MFAI by truncation radius from the city centre.}
\label{tab:t7_truncation}
\begin{tabular}{lrrrrrrrr}
\toprule
City & 10 km & 15 km & 20 km & 25 km & 30 km & 40 km & Full & Swing \\
\midrule
Patna & 0.749000 & 0.748000 & 0.747000 & 0.745000 & 0.746000 & 0.745000 & 0.728000 & 0.021000 \\
Indore & 0.820000 & 0.818000 & 0.817000 & 0.815000 & 0.813000 & 0.810000 & 0.809000 & 0.011000 \\
Varanasi & 0.845000 & 0.838000 & 0.830000 & 0.820000 & 0.816000 & 0.814000 & 0.814000 & 0.031000 \\
Kochi & 0.891000 & 0.885000 & 0.877000 & 0.864000 & 0.854000 & 0.835000 & 0.830000 & 0.061000 \\
Kolkata & 0.898000 & 0.893000 & 0.893000 & 0.893000 & 0.893000 & 0.893000 & 0.893000 & 0.005000 \\
\bottomrule
\end{tabular}
\begin{tablenotes}\footnotesize
\item The ordering is preserved at every radius. Swing is the difference between the 10 km and full-extent values: a large swing indicates that low aggregate accessibility arises from settlement dispersion (Kochi), a small one that it reflects a shortfall of provision (Patna). See Section 6.2.
\end{tablenotes}
\end{table}
